\documentclass[a4paper,12pt]{article}

\title{Blocco 3}
\author{Leonardo Centazzo}
\date{\today}


% Pacchetti essenziali
\usepackage{amsmath, amssymb, amsthm}
\usepackage{enumitem}
\usepackage{geometry}

% Impostazioni pagina minime
\geometry{a4paper, margin=2cm}

% Stili per teoremi, proposizioni, ecc.
\theoremstyle{plain}
\newtheorem{teorema}{Teorema}
\newtheorem{proposizione}{Proposizione}
\newtheorem{lemma}{Lemma}
\newtheorem{corollario}{Corollario}

\theoremstyle{definition}
\newtheorem{definizione}[teorema]{Definizione}
\newtheorem{esempio}[teorema]{Esempio}
\newtheorem{esercizio}{Esercizio}
\newtheorem{osservazione}[teorema]{Osservazione}

\newtheorem*{warning}{Warning}

\theoremstyle{remark}
\newtheorem{nota}[teorema]{Nota}


\begin{document}

\maketitle


\begin{esercizio}
	Let $G$ be any countable group. Consider $2^G, (2^\omega)^G$ equipped with the shift action.
	Prove that there exists an injective, $G$-equivariant function from $2^G$ to $(2^\omega)^G$.
\end{esercizio}

\begin{proof}
	Consider $f: 2^G \to (2^\omega)^G, \ x \mapsto x^{*}$, where
	$\forall g \in G, n < \omega \ x^{*}(g)(n) := x(g)$.
	We want to prove:
	\begin{itemize}
		\item $f$ is $G$-equivariant: \\
			Let $g \in G, x \in 2^{\omega}$. We want to prove that $f(g \cdot x) = g \cdot f(x)$.
			We evaluate both and check they are the same:
			\[
				f(g \cdot x)(h)(n) = (g \cdot x)^{*} (h)(n) =
				(g \cdot x) (h) = x(g^{-1} \cdot h)
			\]

			\[
				(g \cdot f(x))(h)(n) = (g \cdot x^{*})(h)(n) = 
				x^{*}(g^{-1} \cdot h)(n) = x(g^{-1} \cdot n)
			\]

		\item $f$ is injective: \\
			Let $x_1 \not = x_2 \in 2^{\omega}$. Then $x_1 (g) \not = x_2(g)$ for some $g \in G$.
			Then $x_1 ^{*}(g)(0) = x_1 (g) \not = x_2 (g) = x_2 ^{*}(g)(0)$.
			Thus $x_1 ^{*} \not = x_2 ^{*}$.

	\end{itemize}

\end{proof}





\begin{esercizio}
	Let $G,H$ be countable groups and $Y$ a Polish space.
	Prove that if $\{1\} \not = G \leq H$, then:
	\begin{enumerate}
		\item $E(G,Y) \leq_B E(H,Y)$
		\item $F(G,Y) \leq_B F(H,Y)$
	\end{enumerate}
\end{esercizio}

\begin{proof}
	Fix $y_0 \in Y$. Then consider $f: Y^G \to Y^H, \ p \mapsto p^{*}$, where 
	$p^{*}(h) := 
	\begin{cases}
		p(h) & \text{ if } h \in G  \\
		y_0  & else
	\end{cases}
	$

	We must prove:
	\begin{itemize}
		\item $f$ is a reduction from $E(G,Y)$ to $E(H,Y)$ \\
			\begin{itemize}
				\item $\forall p,q \in Y^{G}, g \in G \ \bigg( p = g \cdot q \Rightarrow p^{*} = g \cdot q^{*} \bigg)$ \\

					Let's consider two cases:
					\begin{itemize}
						\item $h \in G$
							\[
								p^{*} (h) = p(h) = (g \cdot q)(h) = 
								q(g^{-1} \cdot h) = q^{*}(g^{-1} \cdot h) =
								(g \cdot q^{*})(h)
							\]
						\item $h \not \in G$ 
							\[
								p^{*} (h) = y_0 = q^{*}(g^{-1} \cdot h) = (g \cdot q^{*})(h)
							\]

					\end{itemize}

				\item $\forall p,q \in Y^{G}, h \in H \ 
					\bigg( p^{*} = h \cdot q^{*} \Rightarrow (\exists g \in G \ p^{*} = g \cdot q^{*}) \bigg)$ \\

					Let $h \in H \setminus G$ such that $p^{*} = h \cdot q^{*}$.
					Then 
					\[
						\forall t \in H \ \ 
						p^{*}(t) = (h \cdot q^{*})(t) = q^{*}(h^{-1} \cdot t)
					\]
					
					In particular 
					\[
						\forall g \in G \ \ y_0 = p^{*}(h \cdot g) = q^{*}(g) = q(g)
					\]

					Then $q = cost_{y_0}$. Then $q^{*} = cost_{y_0}$. Then $p^{*} = cost_{y_0}$.
					Then $p^{*} = 1_{H} \cdot q^{*}$.


				\item $\forall p,q \in Y^{G}, g \in G 
					\bigg( p^{*} = g \cdot q^{*} \Rightarrow p = g \cdot q \bigg) $ 

					\begin{align*}
						p^{*} = g \cdot q^{*} \Rightarrow & \forall h \in H \ p^{*}(h) = q^{*}(g^{-1} \cdot h) \\
						\Rightarrow & \forall t \in G \ p(t) = p^{*}(t) = q^{*}(g^{-1} \cdot t) = 
										q(g^{-1} \cdot t) = (g \cdot q)(t)
					\end{align*}
			\end{itemize}

		\item $f$ is Borel \\
			\[
				(s,r) \in f \iff \bigg( (\forall g \in G \ s(g) = r(g)) \land (\forall h \in H \setminus G \ r(h) = y_0) \bigg)
			\]

	\end{itemize}


\end{proof}





\begin{esercizio}
	Let $E,F$ cber on $X$. Prove that if $E \subseteq F$ and $F$ is treeable, then $E$ is also treeable.
\end{esercizio}

\begin{proof}

\end{proof}





\begin{esercizio}
	Let $X,Y$ be standard Borel spaces, $a:G \curvearrowright X , b:G \curvearrowright Y $ be Borel actions on $X,Y$ respectively.
	Suppose that $f:X \to Y$ be orbit-wise injective, $G$ equivariant, Borel map.
	Prove that if $E_b$ has the cocycle property, then so does $E_a$.
\end{esercizio}

\begin{proof}

\end{proof}


\end{document}

